\subsection{Hyperbola: two complete branches and two asymptotes}

For $a>0$ and $b>0$, the standard hyperbola is
\[
\boxed{\frac{x^2}{a^2}-\frac{y^2}{b^2}=1}.
\]
Define
\[
c=\sqrt{a^2+b^2}.
\]
The foci are $F_1=(-c,0)$ and $F_2=(c,0)$, and the geometric condition is
\[
\boxed{|d(P,F_1)-d(P,F_2)|=2a}.
\]
The absolute value is essential: one sign describes the right branch and the
opposite sign describes the left branch.

\begin{center}
\begin{tikzpicture}[scale=0.82,>=stealth]
    \draw[axis,->] (-6.0,0)--(6.2,0) node[right] {$x$};
    \draw[axis,->] (0,-4.3)--(0,4.5) node[above] {$y$};
    \draw[helper,dashed] (-5.5,-3.67)--(5.5,3.67)
        node[right] {$y=(b/a)x$};
    \draw[helper,dashed] (-5.5,3.67)--(5.5,-3.67)
        node[right] {$y=-(b/a)x$};
    \draw[subset,domain=2.05:5.5,samples=120]
        plot (\x,{1.333*sqrt(\x*\x/4-1)});
    \draw[subset,domain=2.05:5.5,samples=120]
        plot (\x,{-1.333*sqrt(\x*\x/4-1)});
    \draw[subset,domain=-5.5:-2.05,samples=120]
        plot (\x,{1.333*sqrt(\x*\x/4-1)});
    \draw[subset,domain=-5.5:-2.05,samples=120]
        plot (\x,{-1.333*sqrt(\x*\x/4-1)});
    \coordinate (F1) at (-3.33,0);
    \coordinate (F2) at (3.33,0);
    \coordinate (P) at (3.2,1.67);
    \draw[blue!70!black,line width=1pt] (F1)--(P);
    \draw[orange!85!black,line width=1pt] (F2)--(P);
    \fill[geometricSubsetColor] (F1) circle (2.2pt) node[below] {$F_1$};
    \fill[geometricSubsetColor] (F2) circle (2.2pt) node[below] {$F_2$};
    \fill[geometricSubsetColor] (P) circle (2.2pt) node[above right] {$P$};
    \node at (3.8,-3.7) {right branch};
    \node at (-3.8,-3.7) {left branch};
\end{tikzpicture}
\end{center}

The asymptotes
\[
\boxed{y=\pm\frac ba x}
\]
are not parts of the hyperbola. They record the directions approached by both
branches far from the centre.

\begin{corebox}[Do not confuse a branch with the whole hyperbola]
A focus-centred polar equation will naturally describe the branch visible from
that chosen focus and directrix. The Cartesian equation above remains the cleanest
single picture of the complete two-branch curve. Later we will state explicitly
which branch a polar formula traces and how the second branch is recovered.
\end{corebox}
